\documentclass[11pt]{article}
\usepackage[margin=1in]{geometry}
\usepackage[T1]{fontenc}
\usepackage[utf8]{inputenc}
\usepackage{booktabs}
\usepackage{array}
\usepackage{xurl}
\usepackage{xcolor}
\usepackage{hyperref}
\usepackage{enumitem}
\usepackage{lmodern}
\usepackage[expansion=false]{microtype}
\setlist{itemsep=2pt,topsep=4pt}
\emergencystretch=3em
\hypersetup{colorlinks=true,linkcolor=blue!60!black,urlcolor=blue!60!black}
\newcommand{\src}[1]{[#1]}
\newcommand{\HH}{HIV-1}
\newcommand{\nstrains}{24}
\newcommand{\nguides}{190}
\newcommand{\nsp}{150}
\newcommand{\nsa}{40}
\newcommand{\nexact}{0}
\newcommand{\naplus}{63}
\newcommand{\na}{102}
\newcommand{\nb}{25}
\newcommand{\ktwenty}{261}
\newcommand{\kthirty}{0}
\newcommand{\ktx}{9}
\newcommand{\excisebp}{4,439}
\newcommand{\guidesall}{1365}


\title{An In-Silico Pipeline for CRISPR--Cas9 Guide Design Against the \HH{}
  Proviral Genome: Cross-Subtype Conservation, Genome-Wide Human Off-Target
  Exclusion, and Transcriptome Similarity Screening}
\author{Walker, ND\thanks{Independent research; Oregon, USA. Computational
  work performed on a NVIDIA Jetson Orin Nano (ARM64, 8\,GB). All code, data,
  results tables, and the citation ledger are in the accompanying repository
  (\texttt{crispr-hiv-ward}).}}
\date{October 22, 2026}

\begin{document}
\maketitle

\begin{abstract}
Antiretroviral therapy suppresses but does not eradicate HIV-1, and the
first-in-human CRISPR excision trial (EBT-101) was safe but did not provide
durable suppression after treatment interruption. Target and guide-RNA
selection therefore remain open design problems. We built a fully scripted,
reproducible pipeline that (i) star-aligns \nstrains{} diverse HIV-1 complete
genomes to the HXB2 reference to quantify cross-strain conservation of every
candidate guide; (ii) exhaustively screens guides against GRCh38 for
off-target cleavage sites with up to four mismatches (both strands; NGG/NAG
for SpCas9, NNGRRT/NNGRRV for SaCas9); and (iii) measures exact $k$-mer
sharing between the \HH{} genome and both the human genome and the human
protein-coding transcriptome. Of 1,365 enumerated protospacers, 668 passed
design filters and 190 (150 SpCas9, 40 SaCas9) underwent genome-wide
screening. No guide had a single exact-match off-target in the 3.1-billion-base
human genome; 63 guides carried no off-target with $\le$2 mismatches and no
transcriptome hit with $\le$1 mismatch (tier A+). Region-level analysis
identifies the 5$'$ LTR U5 region and the packaging-signal leader as the
combination of high conservation and minimal human similarity, whereas
\emph{pol} and \emph{env} carry substantially more human-shared 20-mers,
consistent with endogenous-retrovirus ancestry. Notably, guides in the
otherwise most conserved leader region exactly match human \emph{NOTCH3}
transcripts and were excluded. We recommend a dual-guide design (5$'$LTR U5 +
integrase; 4,439\,bp excision) with a triple-guide escape-hardened variant,
and provide all sequences, off-target tables, and the evidence ledger.
\end{abstract}

\section{Introduction}

Combination antiretroviral therapy (ART) controls plasma viremia but cannot
reach the latent proviral reservoir, which re-seeds viremia when treatment
stops. Programmable nucleases offered a radical alternative: rather than
blocking replication, excise the proviral DNA itself. Early studies showed
that CRISPR--Cas9 can disrupt or excise integrated \HH{} DNA in cell
lines\src{1}, primary T cells\src{6}, and across the LTR--gag
region\src{2}, and that multiplexed SaCas9 guide RNAs delivered by AAV can
excise proviral fragments in animal models\src{9}. Combining long-acting
ART with CRISPR excision eliminated infectious virus in a subset of
humanized mice\src{4}.

This line of work has reached humans. EBT-101 (Excision BioTherapeutics) is
an AAV9 vector carrying SaCas9 with two guide RNAs targeting the LTR and
gag (construct ``AAV9-SaCas9-GagD-LTR1'')\src{55,20}. Its phase 1/2 trial
(NCT05144386) enrolled six participants, completed in November
2024\src{50,55}, and reported no serious adverse events and no identified
off-target DNA damage\src{18,29}. However, HIV RNA rebounded in all
participants who interrupted ART\src{29,20,21}, and the program is listed as
inactive with the curator note ``safe, but does not provide long-term viral
suppression''\src{55}. The lesson is not that proviral excision failed as a
concept, but that target selection, guide quality, and reservoir coverage
were insufficient\src{22,47}. Two design questions follow directly: which
viral regions should be cut, and how much do candidate guides resemble the
human genome they must not cut?

This paper addresses both questions computationally. We ask: \emph{which
part of the \HH{} genome is the best CRISPR target, considering cross-strain
conservation, functional importance documented in the literature, and least
similarity to human genes?} We contribute: (1) a reproducible pipeline that
quantifies conservation across \nstrains{} diverse complete genomes; (2) an
exhaustive human-genome off-target scan (to four mismatches, both strands,
canonical and non-canonical PAMs) implemented in C and validated guide by
guide; (3) a transcriptome similarity screen that catches human-exact guide
sequences that whole-genome DNA screens under-weight; and (4) a ranked
target-region table and concrete dual/triple guide designs.

\section{Background and related work}

\subsection{Target regions used by published studies}
Published anti-\HH{} CRISPR work has targeted the long terminal repeats
(LTRs)\src{1,6,2}, gag\src{9}, and combinations thereof. Because the
integrated provirus is flanked by two LTRs that are identical copies of the
same sequence, a single LTR-targeting guide cleaves both ends, and NHEJ
re-joining excises the intervening genome\src{9}. Yin et al.\ excised
proviral DNA in vivo with SaCas9 and multiplex gRNAs\src{9}, and Dash et
al.\ estimated excision of approximately 80\% of 5$'$LTR--gag and gag--3$'$LTR
fragments in humanized mice, reporting ``No CRISPR-Cas9 mediated off-target
effects are detected''\src{4}.

\subsection{Viral escape}
Cas9 pressure selects for resistant virus: NHEJ repair of the cut site
inserts indels, and escape variants emerge rapidly when non-conserved
regions are attacked (breakthrough replication in 11--17 days) and more
slowly for conserved targets (20--43 days)\src{32,3}. Resistance mutations
cluster at the cleavage site\src{3}, and host double-strand-break repair
directly generates CRISPR-resistant strains\src{28}. The standard
mitigation is multiplexing several conserved, independent targets\src{3,22,47}.

\subsection{Off-target risk}
Subtype-specific off-target analyses of LTR guide collections found that
149 of 4,114 viable guides (3.60\%) retained predicted human off-target
cleavage, of which only 14 of 326 sites (4.29\%) were exonic\src{11}.
Unbiased assays (GUIDE-seq, CIRCLE-seq, whole-genome sequencing) have
repeatedly failed to detect notable off-target cleavage for anti-\HH{}
guides\src{27}, and clinical monitoring of EBT-101 participants likewise
identified no off-target DNA damage\src{29}. Human endogenous retroviruses
(HERVs) share ancestry with \emph{gag} and \emph{pol}, but an analysis of
predicted off-targets concluded HERV loci are ``not the driving factor for
off-target editing'' for currently used guides\src{27}; our region-level
similarity data (Section~\ref{sec:similarity}) add quantitative context.

\subsection{RNA-targeting alternatives}
Cas13 systems target viral RNA without cutting DNA: LbuCas13a reduced \HH{}
RNA by 50--80\%\src{8}, and RfxCas13d with crRNAs against conserved
\emph{gag}/\emph{pol} regions inhibited replication by more than 90\% in
primary CD4$^+$ T cells\src{44,16}. Combined RNA- and DNA-targeting
ribonucleoproteins are an active frontier\src{54}. Our pipeline designs
Cas9 guides; Cas13 candidates are out of scope but the same
conservation/off-target machinery applies.

\section{Methods}

\subsection{Reference, regions, and strain set}
All coordinates use the HXB2 reference genome (GenBank K03455, 9,719\,bp)
and its GenBank feature annotation\src{48,49}, cross-checked against the
LANL annotated HXB2 map with per-position Shannon entropy\src{51}. We
define 31 regions: LTR subregions (U3/R/U5, both copies), TAR, the
packaging signal $\psi$, \emph{gag} (p17, p24), \emph{pol} (protease, RT,
RNase H, integrase, cPPT), \emph{vif}, \emph{vpr}, \emph{tat}, \emph{rev},
\emph{vpu}, \emph{env} (gp120, RRE, gp41), and \emph{nef}. The $\psi$
coordinates remain literature-derived and unverified; all others come from
the GenBank record or the LANL map. Note that LANL and GenBank label R-region
boundaries within $\pm$2 nucleotides of each other; we use GenBank
throughout.

We downloaded \nstrains{} near-full-length complete \HH{} genomes (8.5--9.9\,kb)
from NCBI via E-utilities, selected for country diversity across Africa,
Europe, Asia, and the Americas, with subtype labels taken verbatim from
record titles where present, and excluding records titled as mutants,
synthetic constructs, or vectors. The NL4-3 proviral genome was extracted
from the pNL4-3 molecular clone record\src{58}. Together with HXB2 this
yields 25 genomes for conservation scoring.

\subsection{Conservation}
Each strain was globally aligned to HXB2 (Biopython PairwiseAligner; match
$+2$, mismatch $-1$, gap open $-4$, extension $-1$), producing a base at
every HXB2 position (gaps/ambiguity coded separately). For each candidate
guide, conservation is the fraction of strains whose aligned sequence
matches the HXB2-derived guide across the full protospacer \emph{and} whose
PAM matches at the constrained PAM positions --- i.e., the fraction of
strains in which the guide site is present unmutated.

\subsection{Guide enumeration and filters}
Protospacers follow the standard convention: the ordered guide sequence is
the 20-nt (SpCas9) or 21-nt (SaCas9) sequence immediately 5$'$ of the PAM on
the PAM-bearing strand. We enumerated every SpCas9 NGG site and every
canonical SaCas9 NNGRRT site on both strands (1,365 sites). Design filters:
GC content 40--65\%, no homopolymer run of four or more identical bases
(this also enforces the U6 promoter poly-T termination rule of standard
guide design\src{38,35}), and no TTT start. Surviving guides (544 SpCas9,
124 SaCas9) were ranked by conservation, then GC near 50\%; the top 150
SpCas9 and top 40 SaCas9 advanced to genome-wide screening.

\subsection{Human off-target scan}
A purpose-built C tool (\texttt{crispr\_scan}; pthreads, 4-bit packed
genome) exhaustively scanned GRCh38 (Ensembl primary assembly, release
113\src{56}) for all sites with Hamming distance $\le$4 to any guide with a
valid PAM: NGG and NAG for SpCas9, NNGRRT and NNGRRV for SaCas9, both
strands, with N bases always counted as mismatches. The tool passes 13 unit
tests including hand-derived strand/PAM vectors, and every screened guide
was confirmed to hit HXB2 exactly at its expected coordinate and strand
(190/190), validating the conventions end-to-end. The full 3.1-Gbp scan of
190 guides completes in approximately five minutes on six ARM cores.

\paragraph{Safety tiers.} A random 20-mer has roughly 30 by-chance
$\le$3-mismatch sites in a 3.1-Gbp genome, so ``zero $\le$3\,mm sites'' is
not a meaningful bar. We therefore tier guides by the sites that matter:
\textbf{A+} --- no genomic off-target with $\le$2 mismatches \emph{and} no
transcriptome hit with $\le$1 mismatch; \textbf{A} --- no genomic off-target
with $\le$1 mismatch and no exact transcriptome hit; \textbf{B} ---
everything else. Cleavage at $\ge$3 mismatches is weak, and sites at 1--2
mismatches are the residual risk band\src{34,27}.

\subsection{Human-similarity maps}
We measured exact $k$-mer sharing between the HXB2 genome and (a) GRCh38
($k=20$ and $k=30$) and (b) GENCODE v44 protein-coding transcripts\src{57}
($k=20$), using the same packed-scanning tool (rolling $k$-mer, both
orientations, N-excluded). For every off-target site with $\le$2 mismatches
(490 sites), we additionally scanned the protein-coding transcriptome at
$\le$1 mismatch to find guide sequences that could cleave human mRNAs.

\section{Results}

\subsection{Conservation across 25 strains}
Per-position identity to HXB2 ranged from 80.2\% to 97.3\% (NL4-3, the
closest subtype-B genome). Table~\ref{tab:regions} summarizes each region's
median guide conservation alongside the independent LANL entropy
values\src{51}: the two measures agree (integrase and cPPT are the most
conserved coding features in both). The 5$'$ LTR U5 region and the R region
score highest among designable regions (median guide conservation 0.72 and
0.84), and the leader/$\psi$ region reaches 0.56 with individual guides at
0.96.

\begin{table}[ht]\centering\small
\caption{HXB2 regions (GenBank\src{48} / LANL map\src{51} coordinates),
median guide conservation across 25 strains, LANL per-region mean Shannon
entropy\src{51} where published, number of tier-A+ guides, and exact
20-mer sharing with the human genome (hu) and protein-coding transcriptome
(tx). $^{\dagger}$literature-derived, unverified.}
\label{tab:regions}
\begin{tabular}{llllrrrr}
\toprule
Region & HXB2 span & Source & Med.\ cons. & LANL ent. & A+ guides & hu\,20-mer & tx\,20-mer \\
\midrule
TAR & 454--513 & LANL map & 0.16 & -- & 0 & 2 & 0 \\
R\_5 & 454--551 & genbank & 0.84 & -- & 1 & 2 & 0 \\
U3\_5 & 1--453 & genbank & 0.00 & -- & 2 & 9 & 0 \\
U5\_5 & 552--634 & genbank & 0.72 & -- & 5 & 1 & 1 \\
psi & 744--802 & lit.$^{\dagger}$ & 0.92 & -- & 0 & 0 & 0 \\
p17 & 791--1186 & LANL map & 0.20 & 0.3664 & 8 & 14 & 0 \\
p24 & 1186--1878 & LANL map & 0.20 & 0.2073 & 12 & 7 & 0 \\
protease & 2253--2549 & LANL map & 0.20 & 0.2002 & 0 & 5 & 0 \\
RT & 2550--3869 & LANL map & 0.12 & 0.2022 & 6 & 53 & 2 \\
RNaseH & 3870--4229 & LANL map & 0.12 & 0.2214 & 0 & 17 & 0 \\
IN & 4230--5093 & LANL map & 0.12 & 0.1598 & 4 & 40 & 2 \\
cPPT & 4781--4785 & LANL map & -- & 0.0112 & 0 &  &  \\
vif & 5041--5619 & genbank & 0.16 & 0.2600 & 4 & 3 & 0 \\
vpr & 5559--5795 & genbank & 0.04 & 0.2616 & 1 & 1 & 0 \\
tat\_e1 & 5831--6045 & genbank & 0.24 & 0.3129 & 2 & 2 & 0 \\
rev\_e1 & 5970--6045 & genbank & -- & 0.3299 & 0 &  &  \\
vpu & 6062--6310 & LANL map & 0.04 & 0.4378 & 0 & 12 & 2 \\
gp120 & 6309--7757 & LANL map & 0.08 & 0.4433 & 2 & 42 & 0 \\
RRE & 7710--8061 & LANL map & 0.12 & 0.2156 & 6 & 8 & 0 \\
gp41 & 7758--8792 & LANL map & 0.04 & 0.3234 & 3 & 31 & 1 \\
nef & 8797--9168 & genbank & 0.04 & -- & 0 & 7 & 0 \\
U3\_3 & 9086--9539 & genbank & 0.04 & -- & 2 & 6 & 0 \\
R\_3 & 9540--9636 & genbank & 0.24 & -- & 1 & 2 & 0 \\
U5\_3 & 9637--9719 & genbank & 0.44 & -- & 3 & 0 & 0 \\
\bottomrule
\end{tabular}

\end{table}

\subsection{Genome-wide off-target screen}
Across all \nguides{} screened guides there were \textbf{zero exact-match
(mm=0) off-target sites} in GRCh38. The full mismatch distribution was:
20 sites at 1 mismatch (17 guides), 470 at 2, 6,773 at 3, and 76,105 at 4.
Tiers: \naplus{} A+, \na{} A, \nb{} B. The single-mismatch sites are listed
per guide in the repository; no guide carries more than two.

\subsection{Human-similarity maps}
\label{sec:similarity}
The \HH{} genome is strikingly alien to human DNA: \textbf{not a single
exact 30-mer is shared} between HXB2 and the 3.1-Gbp human genome, and only
\ktwenty{} of 9,700 possible 20-mer start positions (2.7\%) share an exact
20-mer with it. Sharing is concentrated where retroviral ancestry predicts:
RT (53 human-shared 20-mer positions), gp120 (42), integrase (40), gp41
(31), and p17 (14), versus 1 for 5$'$LTR U5, 0 for $\psi$ and 3$'$LTR U5
(Table~\ref{tab:regions}).

The transcriptome screen proves this matters at guide resolution: among the
490 off-target windows with $\le$2 mismatches, 225 have a $\le$1-mismatch
match in protein-coding transcripts. Critically, the two most conserved
guides in the entire genome --- SP\_772+ and SP\_771+ in the
packaging-signal leader, 0.96 and 0.92 conserved across strains --- carry
\emph{exact} (mm=0) matches to human \emph{NOTCH3} transcripts
(NOTCH3-201) and would have passed a conservation-first design. They were
excluded by the transcriptome screen. Other exact transcript matches
involve \emph{NFX1}, \emph{CHRNA2}, and olfactory-receptor loci.

\subsection{Recommended designs}
\label{sec:strategy}
Table~\ref{tab:strategy} lists the recommended multiplex designs, all
tier A+. The dual design pairs SP\_628+ (5$'$LTR U5, 0.88 conserved, zero
off-targets $\le$2\,mm) with SA\_5067+ (integrase, 0.48) for a 4,439\,bp
excision; because in vivo proviral LTRs are identical, the U5 guide cuts
both LTR ends, and the triple design adds SP\_7794+ (RRE, 0.76) so that
three independent conserved sites must mutate simultaneously for escape,
directly implementing the multiplexing mitigation\src{3,32,22}. Runners-up
for the internal slot were p24 guides (12 tier-A+ guides, 7 human-shared
20-mers, zero transcriptome hits), vif (0.44), and p17 (0.44).

\begin{table}[ht]\centering\small
\caption{Recommended tier-A+ guide combinations (protospacer/PAM as read on
the PAM-bearing strand). Cons.\ = fraction of 25 strains with the guide site
unmutated.}
\label{tab:strategy}
\begin{tabular}{p{2.6cm}p{5.6cm}p{2.2cm}p{1.8cm}}
\toprule
Strategy & Guides (protospacer/PAM) & Excised span & Min.\ cons. \\
\midrule
dual excision & \texttt{TCTAGCAGTGGCGCCCGAAC}--\texttt{AGG} (SP\_628+) \newline \texttt{TGTGTGGCAAGTAGACAGGAT}--\texttt{GAGGAT} (SA\_5067+) & 4,439\,bp & 0.48 \\
triple-guide escape-hardened excision & \texttt{TCTAGCAGTGGCGCCCGAAC}--\texttt{AGG} (SP\_628+) \newline \texttt{TGTGTGGCAAGTAGACAGGAT}--\texttt{GAGGAT} (SA\_5067+) \newline \texttt{GGAGCAGCAGGAAGCACTAT}--\texttt{GGG} (SP\_7794+) & multi-fragment (three cut sites) & 0.48 \\
\bottomrule
\end{tabular}

\end{table}

\begin{table}[ht]\centering\small
\caption{Top tier-A+ guides by conservation. Full tables for all 190 guides,
including every off-target site and its coordinate, are in the repository.}
\label{tab:topguides}
\begin{tabular}{llllrrrr}
\toprule
ID & Region & Position & Protospacer ($5'\to3'$) & PAM & System & Cons. & mm$\le$4 \\
\midrule
SP\_628+ & U5\_5 & 628\,+ & \texttt{TCTAGCAGTGGCGCCCGAAC} & \texttt{AGG} & SpCas9 & 0.88 & 72 \\
SP\_629+ & U5\_5 & 629\,+ & \texttt{CTAGCAGTGGCGCCCGAACA} & \texttt{GGG} & SpCas9 & 0.88 & 73 \\
SP\_7794+ & RRE & 7794\,+ & \texttt{GGAGCAGCAGGAAGCACTAT} & \texttt{GGG} & SpCas9 & 0.76 & 442 \\
SA\_584- & U5\_5 & 584\,- & \texttt{TCCACAGATCAAGGATATCTT} & \texttt{AGATGG} & SaCas9 & 0.72 & 13 \\
SP\_8568+ & gp41 & 8568\,+ & \texttt{GTGGAACTTCTGGGACGCAG} & \texttt{GGG} & SpCas9 & 0.68 & 316 \\
SP\_8567+ & gp41 & 8567\,+ & \texttt{TGTGGAACTTCTGGGACGCA} & \texttt{GGG} & SpCas9 & 0.68 & 337 \\
SP\_796+ & p17 & 796\,+ & \texttt{GCGAGAGCGTCAGTATTAAG} & \texttt{CGG} & SpCas9 & 0.56 & 63 \\
SP\_562+ & U5\_5 & 562\,+ & \texttt{GCCCGTCTGTTGTGTGACTC} & \texttt{TGG} & SpCas9 & 0.52 & 143 \\
SP\_566- & U5\_5 & 566\,- & \texttt{AGCCTTGTGTGTGGTAGATC} & \texttt{TGA} & SpCas9 & 0.52 & 397 \\
SA\_5067+ & IN & 5067\,+ & \texttt{TGTGTGGCAAGTAGACAGGAT} & \texttt{GAGGAT} & SaCas9 & 0.48 & 79 \\
SA\_4382+ & IN & 4382\,+ & \texttt{TGGACAAGTAGACTGTAGTCC} & \texttt{AGGAAT} & SaCas9 & 0.44 & 87 \\
SP\_9647+ & U5\_3 & 9647\,+ & \texttt{GCCCGTCTGTTGTGTGACTC} & \texttt{TGG} & SpCas9 & 0.44 & 143 \\
SP\_5199+ & vif & 5199\,+ & \texttt{AGAAGTACACATCCCACTAG} & \texttt{GGG} & SpCas9 & 0.44 & 204 \\
SP\_884+ & p17 & 884\,+ & \texttt{AACATATAGTATGGGCAAGC} & \texttt{AGG} & SpCas9 & 0.44 & 220 \\
\bottomrule
\end{tabular}

\end{table}

\section{Discussion}

\paragraph{The sweet spot is the LTR U5 region, with p24 as the best
internal partner.} Our data resolve the design tension cleanly. The most
\emph{conserved} guides sit in the leader/$\psi$ region (0.96) --- but the
most conserved guides are exactly the ones that match human transcripts
(\emph{NOTCH3}), a trap for conservation-first pipelines. The 5$'$ LTR U5
region offers near-top conservation (0.72--0.88) with essentially no human
similarity (one human-shared 20-mer position genome-wide), the dual-cut
excision geometry of the LTRs, and five tier-A+ guides. Among internal
regions, p24 combines the largest tier-A+ pool (12) with low human sharing,
whereas \emph{pol} and \emph{env} --- despite excellent functional
essentiality --- carry the most human-shared sequence, as expected from
HERV ancestry; this did not produce off-targets for our specific guides
(consistent with\src{27}), but it lowers the margin of safety for guides
designed naively in those regions.

\paragraph{Escape remains the binding constraint.} EBT-101's rebound after
treatment interruption\src{29,20,21} and the documented speed of NHEJ-driven
escape\src{32,3} both point to the same conclusion: two guides are likely
insufficient, and reservoir coverage --- not guide safety --- is the hard
problem. Proviral DNA circles with reformed LTRs can persist after
excision\src{13}, defective proviruses accumulate from acute infection
onward\src{15}, and intact inducible proviruses are the true target
population\src{39}. Our triple design (LTR U5 + integrase + RRE) maximizes
simultaneous-escape requirements while keeping every guide in tier A+;
whether three guides suffice against the reservoir is an empirical question
that only animal and clinical work can answer\src{22,47,43}.

\paragraph{Relationship to EBT-101.} EBT-101 targets GagD and LTR1 with
SaCas9\src{55,20}. Our pipeline's top LTR guide differs (U5 rather than
their LTR site) because we required zero $\le$2\,mm off-targets and
transcriptome cleanliness; whether the published GagD/LTR1 guides would
pass our tier system is directly testable with the repository pipeline and
is planned as a benchmark. The clinical experience to date --- safe, no
off-target damage, but viral rebound\src{18,29,50} --- is consistent with
our modeling that guide safety is achievable and that reservoir coverage
and escape are the unsolved problems.

\section{Limitations}
This is a computational design study; nothing here is experimentally
validated, and nothing here is medical advice. Specific limitations:
(1) off-target scanning covers mismatches only, not DNA or RNA bulges, and
uses NGG/NAG and NNGRRT/NNGRRV PAM sets --- CHOPCHOP/CRISPOR-grade
re-scoring with bulge models is advisable before any wet-lab use;
(2) conservation uses \nstrains{} near-full-length genomes selected for
geographic spread, not a curated subtype-balanced panel, and subtype labels
come from record titles (only two records carried explicit subtype labels);
(3) the transcriptome screen covers protein-coding transcripts only
(GENCODE v44); non-coding RNAs are unscreened;
(4) $\psi$ coordinates are literature-derived and unverified; LANL and
GenBank differ by $\pm$2 nt on some LTR feature boundaries;
(5) SaCas9's NNGRRT PAM yields few designable sites (124 passing), limiting
AAV-compatible options relative to SpCas9;
(6) no chromatin, nucleosome, or delivery context is modeled; and
(7) viral escape dynamics, reservoir size, and excision efficiency are not
modeled. In silico tier A+ status is a necessary, not sufficient, condition
for a therapeutic guide.

\section{Data and code availability}
The repository (\texttt{crispr-hiv-ward}) contains the full pipeline
(\texttt{src/}), unit tests (39 passing across the Python pipeline and the C
scanner), all result tables (\texttt{results/}), the strain sequences,
packed-genome tooling, the citation ledger with verbatim source quotes
(\texttt{docs/citations/}), and this paper's \LaTeX{} source. The human
genome (Ensembl release 113\src{56}) and GENCODE v44\src{57} are fetched by
script; strain sets are fetched from NCBI E-utilities with reproducible
queries.

\section*{Sources}
\small
Entries below were registered at retrieval time and verified against the fetched page text; quoted passages are verbatim substrings of the saved page files in \texttt{docs/citations/pages/}. Sources without a quoted passage are data repositories or pages that yielded no clean quotable sentence; they are cited from their fetched content. Very long machine-extracted passages (over 200 characters, e.g.\ raw sequence translations) are omitted from the rendered quotes for typography; they remain in the ledger file.

\begin{enumerate}
\item [1] Ebina et al. 2013, Sci Rep: Harnessing the CRISPR/Cas9 system to disrupt latent HIV-1 provirus \\ \url{https://www.nature.com/articles/srep02510}
  \\ {\sloppy \emph{Quote:} ``two gRNA-expressing plasmids were generated for targets 5 and 6 (T5 and T6), located in the TAR sequence of the R region and NF-$\kappa$B binding sequence in the U3 region, respectively''}

\item [2] Nat Commun 2015 (Liao?): Use of the CRISPR/Cas9 system as an intracellular defense against HIV-1 infection in human cells \\ \url{https://www.nature.com/articles/ncomms7413}
  \\ {\sloppy \emph{Quote:} ``Target gRNAs were designed against the enhanced green fluorescent protein (EGFP) coding region as well as the non-coding long terminal repeat (LTR) of an EGFP reporter lentivirus''}

\item [3] Wang et al. 2016, Retrovirology: CRISPR/Cas9: a double-edged sword when used to combat HIV infection \\ \url{https://pmc.ncbi.nlm.nih.gov/articles/PMC4882869}
  \\ {\sloppy \emph{Quote:} ``The majority of the resistance mutations appeared to cluster at the site at which Cas9 was designed to cleave the viral DNA, even though the sgRNA binding site is much bigger.''}
  \\ {\sloppy \emph{Quote:} ``Another striking feature was the frequent occurrence of insertions and deletions (indels), at least for the less conserved viral target sequences.''}

\item [4] Dash et al. 2019, Nat Commun: Sequential LASER ART and CRISPR treatments eliminate HIV-1 in a subset of infected humanized mice \\ \url{https://www.nature.com/articles/s41467-019-10366-y}
  \\ {\sloppy \emph{Quote:} ``No CRISPR-Cas9 mediated off-target effects are detected.''}
  \\ {\sloppy \emph{Quote:} ``Excision efficiency was estimated to be 80\% in both 5$'$-LTR-Gag and Gag to 3'-LTR in mouse M4349.''}

\item [6] Hu et al. 2014, PNAS: RNA-directed gene editing specifically eradicates latent and prevents new HIV-1 infection \\ \url{https://www.pnas.org/doi/10.1073/pnas.1405186111}
  \\ {\sloppy \emph{Quote:} ``genotoxicity nor off-target editing to the host cells, and completely excised a 9,709-bp fragment of integrated proviral DNA that spanned from its 5$'$ to 3$'$ LTRs.''}

\item [8] CRISPR-Cas13a Inhibits HIV-1 Infection \\ \url{https://pmc.ncbi.nlm.nih.gov/articles/PMC7321785}

\item [9] Yin et al. 2017, Mol Ther: In Vivo Excision of HIV-1 Provirus by saCas9 and Multiplex Single-Guide RNAs in Animal Models \\ \url{https://pmc.ncbi.nlm.nih.gov/articles/PMC5417847}

\item [11] Prediction of HIV-1 Subtype-Specific Off-Target Effects (J Virol 2019) \\ \url{https://pmc.ncbi.nlm.nih.gov/articles/PMC6553478}
  \\ {\sloppy \emph{Quote:} ``NAs with predicted off-target cleavage was calculated to be 149 of the 4,114 viable gRNAs that were created, which corresponds to 3.60\% of the tested gRNAs.''}
  \\ {\sloppy \emph{Quote:} ``Of the 326 off-target effects, only 14 (4.29\%) of them were exonic''}
  \\ {\sloppy \emph{Quote:} ``gRNA templates were designed from every possible 20-mer from a collection of patient HIV-1 5$'$ LTR sequences.''}
  \\ {\sloppy \emph{Quote:} ``sidered viable for a CRISPR-Cas9 treatment, this means that 96.40\% of the viable gRNAs are predicted to generate no off-target effects''}

\item [13] J Virol 2021: CRISPR/Cas9 Ablation of Integrated HIV-1 Accumulates Proviral DNA Circles with Reforming LTRs (PMC) \\ \url{https://pmc.ncbi.nlm.nih.gov/articles/PMC8577360}
  \\ {\sloppy \emph{Quote:} ``The T5 gRNA targets a highly conserved region between the TATA box and the NF-$\kappa$B binding site in LTRs.''}
  \\ {\sloppy \emph{Quote:} ``geting CRISPR/Cas9, we show that the excised HIV-1 provirus persists for a few weeks and may rearrange in circular molecules.''}

\item [15] Defective proviruses rapidly accumulate during acute HIV-1 infection (PMC5014606) \\ \url{https://www.ncbi.nlm.nih.gov/pmc/articles/PMC5014606}
  \\ {\sloppy \emph{Quote:} ``cumulate rapidly within the first few weeks of infection to make up over 93\% of all proviruses, regardless of how early ART is initiated.''}
  \\ {\sloppy \emph{Quote:} ``The most common defects were internal deletions (80\%), which varied in size (15 bp to \textasciitilde{}8 kb) and in location in the genome.''}
  \\ {\sloppy \emph{Quote:} ``Some proviruses had 3$'$ deletions affecting the \_env\_, \_tat\_, \_rev\_ and \_nef\_ genes.''}

\item [27] Frontiers review 2020: Designing Safer CRISPR/Cas9 Therapeutics for HIV \\ \url{https://pmc.ncbi.nlm.nih.gov/articles/PMC7434968}
  \\ {\sloppy \emph{Quote:} ``ed with the human genome and have a high similarity to HIV, they are not the driving factor for off-target editing.''}

\item [28] Host Double Strand Break Repair Generates HIV-1 Strains Resistant to CRISPR/Cas9 \\ \url{https://pmc.ncbi.nlm.nih.gov/articles/PMC4941621}
  \\ {\sloppy \emph{Quote:} ``ions did lead to HIV-1 strains resistant to CRISPR editing, many of the resistant strains were due to DSB repair associated mutagenesis.''}

\item [29] Medscape 2024: First-in-Human Trial of CRISPR Gene Therapy for HIV (EBT-101) \\ \url{https://www.medscape.com/viewarticle/first-human-trial-crispr-gene-therapy-hiv-2024a1000dx5}
  \\ {\sloppy \emph{Quote:} ``SUGGESTED FOR YOU "There was no off-target DNA damage that's been identified so far," Presti said.''}
  \\ {\sloppy \emph{Quote:} ``who had antiretroviral interruption, but one patient had a delayed rebound of almost 16 weeks along with a significant drop in the HIV reservoir.''}

\item [39] Elimination of HIV Reservoirs Harboring Intact Proviruses (2025/2026) \\ \url{https://pmc.ncbi.nlm.nih.gov/articles/PMC12746608}

\item [48] GenBank K03455.1: HIV-1 HXB2CG complete genome (NCBI, HXB2 reference sequence) \\ \url{https://www.ncbi.nlm.nih.gov/nuccore/K03455}
  \\ {\sloppy \emph{Quote:} ``LOCUS HIVHXB2CG 9719 bp ss-RNA linear VRL 21-OCT-2002 DEFINITION Human immunodeficiency virus type 1 (HXB2), complete genome; HIV1/HTLV-III/LAV reference genome.''}

\item [51] LANL HIV Database: HXB2 reference sequence (K03455) annotated map with per-position entropy (HXB2.xlsx) \\ \url{https://www.hiv.lanl.gov/content/sequence/HIV/MAP/HXB2.xlsx}
  \\ {\sloppy \emph{Quote:} ``protein (frame 1) | protein (frame 2) | protein (frame 3) | entropy (LANL's 2022 DNA filtered web alignment) | RNA/DNA features | RNA/DNA features notes ...''}

\item [54] Banik et al. 2026 Front Genome Ed: Combined RNA- and DNA-targeting CRISPR ribonucleoproteins enhance suppression of HBV transcripts \\ \url{https://pmc.ncbi.nlm.nih.gov/articles/PMC13642003}

\item [55] SCGE Platform: EBT-101 trial report (AAV9-SaCas9-GagD-LTR1) \\ \url{https://scge.mcw.edu/platform/data/report/clinicalTrials/NCT05144386}
  \\ {\sloppy \emph{Quote:} ``Drug Product Description | AAV9-SaCas9-GagD-LTR1''}
  \\ {\sloppy \emph{Quote:} ``Safe, but does not provide long-term viral suppression''}
  \\ {\sloppy \emph{Quote:} ``Completion Date | 2024-11-14''}

\item [56] Ensembl release 113: Homo sapiens GRCh38 primary assembly DNA (data source) \\ \url{https://ftp.ensembl.org/pub/release-113/fasta/homo_sapiens/dna}

\item [57] GENCODE release 44: protein-coding transcripts (data source) \\ \url{https://ftp.ebi.ac.uk/pub/databases/gencode/Gencode_human/release_44}

\item [58] GenBank M19921: pNL4-3 molecular clone record (NL4-3 proviral genome extracted) \\ \url{https://www.ncbi.nlm.nih.gov/nuccore/M19921}

\end{enumerate}


\end{document}